\section{The arithmetic frame with a bound chosen before depth}
\label{sec:arithmetic-frame-uniform-depth}

The constant in the curve argument is chosen before the
cyclotomic depth.  Finiteness of cohomology at one varying
number field would not justify this order.
We prove the required uniformity and the two strict
classes on the same fields as the seed.

Fix an odd prime \(\ell\), a finite field \(k\) of
characteristic \(\ell\), an imaginary quadratic character
\(\eta\), and a continuous finite-order character
\(\Psi:G_{\mathbb Q}\to k^\times\) with \(\Psi(c)=1\).
Its order is prime to \(\ell\).
In the principal application
\(\eta=\chi_{\mathbb Q(\sqrt{-3})}\) and
\(\Psi=\bar\varepsilon^a\bar\psi
      =\bar\varepsilon^j\bar\chi'\).
The fixed totally real abelian field \(F_b\) has
conductor prime to \(\ell\).
For the principal generic packet it is the
biquadratic field constructed in
Section~\ref{sec:seed-all-generic-primes}.
The argument also applies to a fixed real quadratic
base after the local conditions below have been checked.
It has the following checked properties:
\begin{enumerate}
\item \(F_bK/F_b\), where \(K=\mathbb Q^{\ker\eta}\),
is unramified at every finite place and split at the
prescribed level primes, including \(\ell\).
\item At \(\ell\), \(F_b/\mathbb Q\) is unramified,
\(\Psi|_{I_\ell}=\bar\varepsilon^j\), and at every
place \(v\mid\ell\) of \(F_b\) the character
\(\Psi^2|_{G_{F_{b,v}}}\) is nontrivial.
The generic packet has the stronger property
\(2j\not\equiv0\pmod{\ell-1}\), which verifies this
condition on inertia itself.
For an exceptional coefficient prime the stated
condition may instead be verified on the full
decomposition group.
\item There is a prime \(p_0\), split in \(F_bK\), with
\(\Psi(\operatorname{Frob}_{p_0})=1\) and
\(\bar\varepsilon(\operatorname{Frob}_{p_0})=-1\).
\end{enumerate}
All finite ramification of \(\eta,\Psi\) is included in
a fixed integer \(M\), divisible by \(2\ell\).

Choose the fresh prime \(w\) used for the Eisenstein
seed, with \(w\equiv1\pmod{\ell}\) and with
\(\eta,\Psi,\bar\varepsilon\) trivial on its decomposition
group.  Let \(F_\theta/\mathbb Q\) be its cyclic
degree-\(\ell\) field.
Put
\[
 \begin{gathered}
 C_r=\mathbb Q(\zeta_{\ell^{r+1}})^{\mu_{\ell-1}},
 \qquad F_r^+=F_bF_\theta C_r,\qquad L_r=F_r^+K,\\
 E_0=F_bF_\theta K\,\mathbb Q^{\ker\Psi}\,
                                  \mathbb Q(\mu_\ell),
 \qquad E_r=E_0C_r.
 \end{gathered}                                             \tag{AF1}
\]
The finite abelian field \(E_0\) is fixed before \(r\).
The degrees of \(E_r/F_r^+\) and \(E_r/L_r\) are prime to
\(\ell\).

\subsection{One class-group bound throughout the cyclotomic tower}
\begin{srlemma}\label{lem:frame-class-rank-uniform}
There is \(C_0\ge0\), depending on the fixed field \(E_0\)
and on \(\ell\), such that
\[
       \dim_{\mathbb F_\ell}
             \operatorname{Cl}(E_r)/\ell\operatorname{Cl}(E_r)
                                      \le C_0
                         \qquad(r\ge0).                       \tag{AF2}
\]
No bound on the numerical size of \(C_0\) is needed here.
\end{srlemma}
\begin{proof}
We use the classical cyclotomic Iwasawa module and the
\href{https://doi.org/10.2307/1971116}
{Ferrero--Washington theorem} in its actual scope:
the base field \(E_0\) is a finite abelian extension of
\(\mathbb Q\), and the tower is its cyclotomic
\(\mathbb Z_\ell\)-extension.
Write
\[
 A_r=\operatorname{Cl}(E_r)[\ell^\infty],
       \qquad X=\varprojlim_r A_r
\]
with transition maps given by norms.
The Iwasawa module theorem makes \(X\) a finitely
generated torsion \(\mathbb Z_\ell[[T]]\)-module.
Ferrero--Washington gives \(\mu(X)=0\).
The structure theorem therefore makes \(X\) a finite
\(\mathbb Z_\ell\)-module in the sense of finite generation.
Indeed the elementary modules
\(\mathbb Z_\ell[[T]]/(f(T)^e)\), with \(f\) distinguished,
are finite free over \(\mathbb Z_\ell\), and the finite
kernels and cokernels in the structure theorem preserve
finite generation.

For completeness, projections to finite class groups
are surjective after discarding finitely many initial
layers.
In a cyclotomic tower all primes above \(\ell\) are
totally ramified beyond a fixed layer.
Thus no nontrivial subextension of \(E_t/E_r\), for
\(t\ge r\) beyond that layer, is unramified.
Its intersection with the Hilbert class field of \(E_r\)
is \(E_r\).
Given any ideal class of \(E_r\), Chebotarev supplies a
prime in that class which splits in \(E_t/E_r\).
A prime above it has that class as its norm.
The norm on class groups is consequently surjective,
and so is the norm on their \(\ell\)-primary parts.
Compactness of the finite inverse system now gives
\(X\twoheadrightarrow A_r\).
Hence
\(\dim_{\mathbb F_\ell}A_r/\ell A_r
 \le\dim_{\mathbb F_\ell}X/\ell X\)
for all those layers.
Enlarge this bound to cover the finitely many initial
layers.  This proves (AF2).
\end{proof}

This is the actual argument of
\srcid{FrobeniusSmoothVectors.RankLeOfCardLe.lev\_rk}
and \srcid{b0}.
The source fixes
\[
 V=H_0\cap H_{\rm th}(w)\cap H_b\cap
                         \ker\eta\cap\ker\Psi\cap\ker\bar\varepsilon
\]
and uses the exact equality
\[
 U_r\cap\ker\eta\cap\ker\Psi\cap\ker\bar\varepsilon
              =V\cap\ker(\varepsilon\bmod\ell^{r+1}).           \tag{AF3}
\]
The fixed field of \(V\) is the fixed abelian base in
that proof.  Formula (AF1) permits an additional fixed
real abelian factor; the same proof applies verbatim.
A tower of unrelated fresh odd conductors would not
give (AF2) from Ferrero--Washington.

\subsection{The three cohomology bounds defining the frame}
Let \(\Gamma_r\) be the quotient of \(G_{F_r^+}\) by the
closed normal subgroup generated, away from \(\ell\), by
the inertia elements on which both \(\eta\) and \(\Psi\)
are trivial.
Let \(\delta\) denote its induced quadratic character,
\(\Gamma_{r,0}=\ker\delta\), and use the actual local
decomposition groups above \(\ell\) as local maps.
For the CM subgroup use both members of each conjugate
pair of places.

For a representation \(W\) and stable local subspaces
\(b_v\), denote by \(H_C^i(W,b)\) the mapping-cone
cohomology of
\[
 C^\bullet(\Gamma_r,W)\longrightarrow
          \prod_{v\mid\ell}C^\bullet(G_{F_{r,v}^+},W/b_v),
\]
with the cone shifted so that its degree-one elements
are cocycles \(z\) together with local correcting
vectors \(Y_v\), modulo global coboundaries.
When \(b_v=0\) it is the strict local condition.

\begin{srlemma}\label{lem:frame-three-bounds}
After enlarging the same \(C_0\) if necessary,
simultaneously for every depth \(r\),
\[
 \begin{split}
 \dim_kH_C^1(\Gamma_r,k(\Psi^{-2}\delta),0)&\le C_0,\\
 \dim_kH_C^1(\Gamma_{r,0},k(\Psi^{-1}),0)&\le C_0,\\
 \dim_kH^1(\Gamma_r,k)&\le1+C_0 .
 \end{split}                                                 \tag{AF4}
\]
These are precisely \srcid{Frame.Ok.b1p},
\srcid{b1}, and \srcid{h1}.
\end{srlemma}
\begin{proof}
At every ordinary local group, \(\Psi^2\ne1\).
Indeed this was checked over \(F_b\), and the subsequent
local extensions have \(\ell\)-power degree; restriction
through them cannot decrease the nontrivial image of a
character of order prime to \(\ell\).
The local invariant spaces in the first two lines are
therefore zero.
Consequently forgetting the local correcting vectors
injects their mapping-cone \(H^1\) into global \(H^1\):
two local primitives of a global boundary differ by a
local invariant, which is zero.

Restrict further to \(G_{E_r}\).
The intermediate quotient has order prime to \(\ell\),
so averaging kills the kernel in inflation--restriction.
Both coefficient characters are now trivial.
The quotient defining \(\Gamma_r\) kills inertia away
from \(\ell\) once \(\eta,\Psi\) are trivial, and the
strict condition kills inertia above \(\ell\).
Thus the restrictions are everywhere-unramified
homomorphisms \(G_{E_r}\to k\).
Class field theory identifies them with subspaces of
\(\operatorname{Hom}(\operatorname{Cl}(E_r),k)\).
Its dimension over \(k\) is the
\(\mathbb F_\ell\)-dimension in (AF2).
This proves the first two bounds.

For the third line, restriction injects into
\(H^1(G_{E_r,S_\ell},k)\).
Kummer theory, followed by the finite prime-to-\(\ell\)
quotient invariants, identifies its possible image
with the cyclotomic-character part of
\[
             E_r^\times/(E_r^\times)^\ell
\]
which is unramified away from \(\ell\).
There is the exact Kummer sequence with terms
\[
 \mathcal O_{E_r}[1/\ell]^\times/
                       (\mathcal O_{E_r}[1/\ell]^\times)^\ell,
 \qquad
 \operatorname{Cl}(\mathcal O_{E_r}[1/\ell])[\ell].
                                                               \tag{AF5}
\]
The class-group term has dimension at most \(C_0\).

The cyclotomic part of the valuation module at
\(\ell\) is zero.
Indeed the tame cyclotomic inertia is nontrivial at
each such place over \(F_r^+\); it fixes that place
but acts nontrivially through \(\bar\varepsilon\).
The coefficient of an eigenvector at that place must
therefore vanish.
Hence the relevant \(S_\ell\)-unit part in (AF5)
comes from ordinary units.

The field \(E_r\) is CM.
For any unit \(u\), the unit \(u/\overline u\) has
absolute value one at every archimedean embedding;
it is a root of unity.
Since two is invertible modulo \(\ell\), the
minus part of units modulo \(\ell\) comes from
roots of unity.
The character \(\bar\varepsilon\) is odd, and the
\(\ell\)-primary roots of unity form a cyclic group.
Its contribution has dimension at most one.
This proves the final bound in (AF4).
\end{proof}

\subsection{The two actual nonzero strict classes}
The two nontriviality fields of the frame do not require
assuming the existence of an unexplained deformation.
They are forced by one additional ramified place.

\begin{srlemma}\label{lem:frame-strict-classes}
For \(a'=1,2\), put \(\kappa_{a'}=\Psi^{-a'}\eta\).
The quotient \(\Gamma_r\) has
\[
                H_C^1(\Gamma_r,k(\kappa_{a'}),0)\ne0
                                  \qquad(r\ge0).              \tag{AF6}
\]
\end{srlemma}
\begin{proof}
The character \(\kappa_{a'}\) is odd.
It is nontrivial on every full local decomposition
group at \(\ell\), by the displayed \(\Psi^2\ne1\)
condition and splitting of the CM extension there.
Nontriviality on inertia is not required.
It differs from \(\bar\varepsilon\): at the selected
scalar prime \(p_0\) it has value one, whereas
\(\bar\varepsilon\) has value minus one.

Over \(\mathbb Q\), take the finite ramification set
consisting of \(\ell,w\), and the primes ramifying in
\(\kappa_{a'}\).
Use zero local conditions at \(\ell\) and at the
ramified primes, and the full local \(H^1\) at \(w\).
The Selmer dimension formula from finite
Poitou--Tate duality gives
\[
 \dim_k H^1_{\mathcal L}(\mathbb Q,k(\kappa_{a'}))
 -\dim_k H^1_{\mathcal L^\perp}
           (\mathbb Q,k(\kappa_{a'}^{-1}\bar\varepsilon))
                                                        =1 . \tag{AF7}
\]
Here are all correction terms.
Both global invariant spaces are zero because
\(\kappa_{a'}\ne1,\bar\varepsilon\).
The archimedean invariant space is zero because
\(\kappa_{a'}\) is odd.
At \(\ell\), and at the primes ramifying in
\(\kappa_{a'}\), local invariants are zero.
At \(w\), both \(\kappa_{a'}\) and
\(\bar\varepsilon\) are trivial; local duality and the
local Euler formula give \(h^1=2\), \(h^0=1\).
Thus \(w\) contributes exactly one.
At unramified primes outside the set, the unramified
condition has dimension \(h^0\) and contributes zero.
These are exactly the terms in (AF7).
In particular there is a nonzero Selmer class \(z\).

The class stays nonzero on \(G_{F_r^+}\).
The extension \(F_r^+/\mathbb Q\) is Galois, and
\(\kappa_{a'}\) is nontrivial on a full ordinary local
decomposition group over \(F_r^+\), as just proved.
Its invariant space over \(G_{F_r^+}\) is therefore zero.
Inflation--restriction therefore makes restriction
on \(H^1\) injective.

At \(w\), the character is trivial, and a cocycle
restricts on inertia to a homomorphism into the
additive group of \(k\).
Wild inertia has pro-\(w\) image and is killed.
The remaining \(\ell\)-part of tame inertia is
\(\mathbb Z_\ell(1)\).
The degree-\(\ell\) field \(F_\theta\) is totally
ramified at \(w\); inertia after restriction to it
maps into \(\ell\mathbb Z_\ell(1)\).
Our homomorphism is zero there.
At every other finite prime different from \(\ell\),
a local boundary vanishes on inertia whenever its
coefficient character is one; at primes outside
the ramification set the class is unramified.
Thus \(z\) vanishes on exactly the inertia elements
normally killed in defining \(\Gamma_r\).
Products, conjugates and continuity give vanishing
on their closed normal closure.
The class descends to \(\Gamma_r\).
It is locally a boundary at every place above
\(\ell\), so it supplies (AF6), including its actual
local correcting vectors.
No nonzero target character group, large Selmer
length, or desired character estimate was assumed
in this construction.
\end{proof}

\subsection{Completing the frame and then choosing depth}
We record the remaining arithmetic assertions, to
make the use of \srcid{Frame.Ok} explicit.

The group \(\Gamma_r\) is profinite.
Its quadratic character \(\delta\) has open kernel,
complex conjugation is an involution with
\(\delta(c)=-1\), and the induced characters
\(\Psi,\varepsilon\) are continuous.
There are \(d=[F_r^+:\mathbb Q]\) archimedean places,
and the sum of the ordinary local degrees is \(d\).
The marked element \(g_0\) with \(\Psi(g_0)^2\ne1\)
may be selected in a full ordinary local decomposition
group at \(\ell\).  In the generic packet it may be
selected in inertia.

Every residue degree at \(p_0\) in \(F_r^+K\) is odd:
the fixed real quadratic factors split there, and
the remaining extensions have \(\ell\)-power degree.
Thus a local Frobenius \(s_0\in\ker\delta\) has
\[
                         \Psi(s_0)=1,\qquad
                  \bar\varepsilon(s_0)=-1 .                  \tag{AF8}
\]
It proves all three nonvanishing requirements
\(\bar\varepsilon\ne1\),
\(\bar\varepsilon\Psi^{\pm1}\ne1\), and
\(\bar\varepsilon\Psi^{\pm2}\delta\ne1\).
It also excludes
\(\bar\varepsilon=\Psi^b\) on \(\ker\delta\)
for every integer \(b\).
Local square separation holds at every ordinary place
by the full-decomposition condition above and its
preservation under the \(\ell\)-power layers.
On the inertia kernel of \(\varepsilon\), \(\Psi=1\).

A completion of \(F_r^+\) at \(\ell\) contains no
\(\mu_\ell\): its ramification index is an
\(\ell\)-power.
Its pro-\(\ell\) unit group is therefore a free
\(\mathbb Z_\ell\)-module of rank \(d_v\).
Choose a basis adapted to the cyclotomic norm,
whose kernel has rank \(d_v-1\).
All basis elements lie in the kernel of \(\Psi\),
and all but one lie in the kernel of \(\varepsilon\).
Any continuous character into a finite abelian
\(\ell\)-group vanishing on this basis vanishes on
inertia.
This is the exact \srcid{Frame.Ok.gen} assertion.
A pro-\(\ell\) character of \(\ker\delta\) unramified
at every ordinary place is unramified at all finite
places: the other allowed inertia quotients have
order prime to \(\ell\).
It factors through the finite class group of \(L_r\).
Its exponent supplies \srcid{Frame.Ok.cl}.

The finite-cohomology condition \(\Phi_\ell\) follows
from the ordinary arithmetic finiteness theorem on
every open subgroup: its finite representations
are unramified outside a finite set, and restriction
to a finite splitting field reduces degree-one
finiteness to the unit and class-group calculation.
This is a finiteness assertion, not a uniform
bound in \(r\).

Finally the Euler--Poitou--Tate comparison for the
same global quotient and these local maps gives,
for a finite \(k'\)-representation \(W\) and stable
local subspaces \(b_v\), with
\(\operatorname{Hom}_{\Gamma_r}
       (W,k'(\bar\varepsilon))=0\),
\[
 h_C^2(W,b)+h_C^0(W,b)+\sum_v d_v\dim b_v
       \le h_C^1(W,b)+\sum_{u\mid\infty}\dim W^{c_u}.            \tag{AF9}
\]
This is the precise \srcid{Frame.Ok.ec} statement.
It uses the standard finite global and local duality
already proved in the cohomology part of the manuscript.
The local unramified conditions away from \(\ell\)
have Euler contribution zero.
The last possible cone degree is the global dual
invariant term, which vanishes by the displayed
\(\operatorname{Hom}\) hypothesis.
Thus the degree-two cone is finite and its Euler
formula gives (AF9).
For the normalized residual matrix representations
used below, the dual invariant hypothesis is checked
by triangular entries: the diagonal differences are
\(0,\pm1,\pm2\); (AF8) excludes the first two sorts,
and complex conjugation gives the required
\(\delta\)-factor on the opposite-corner
\(\pm2\) entries.
This is the matrix proof
\srcid{FrobeniusSmoothVectors.PrDual.pr\_dual}.

Equations (AF4), (AF6), and the preceding assertions
verify every field of \srcid{Frame.Ok}\((C_0)\).
Only now choose \(r\) with
\[
                        \ell^r>3C_0+7 .                       \tag{AF10}
\]
Since every \(d_v\ge\ell^r\) and \(d\ge d_v\),
\[
                         d>3C_0+7,\qquad d_v>5 .
\]
These are exactly \srcid{Frame.Big}\((C_0)\).
Any additional prescribed global or local depth
bound may be imposed by further increasing \(r\);
the same \(C_0\) still applies.

\paragraph{Exact source correspondence and scope.}
The actual producers are
\srcid{RankLeOfCardLe.b0},
\srcid{FrameEx.ok\_link},
\srcid{FrameEx.frame\_ex},
\srcid{Strict.st\_char}, \srcid{Strict.strict},
and
\srcid{FrobeniusGaloisFieldFourCompatibility}.
They are under
\srcid{Development.AutomorphicArithmetic.AutomorphicVarieties.%
FrobeniusSmoothVectors}
at the frozen tag \srcid{alpoge-source-7cbf7a0c}.
The three comparison bounds are assembled by
\srcid{GaloisCochains.B1p},
\srcid{B1}, and \srcid{H1}.
The additional fixed real abelian factors in (AF1)
do not require a new Iwasawa theorem:
their compositum is absorbed into \(E_0\) before
the cyclotomic depth is chosen.

This proves the actual arithmetic frame used by
the curve and sharp-ideal arguments on the
principal nonresonant packet, and on any exceptional
packet for which the same displayed fixed-base and
full-decomposition conditions have been verified.
It supplies auxiliary existence and depth; it
does not assert a polynomial numerical bound for
\(C_0\), nor use such a bound in the final
host-uniform character estimate.
